% Chapter 3

\chapter{Methodology} % Main chapter title

\label{Chapter3}

%----------------------------------------------------------------------------------------

\section{Research Design}

The study follows design science research (DSR), which creates and evaluates an artefact intended to solve an identified
organisational problem \parencite{hevner2004design}. It is organised in the six activities of the DSR methodology of
\textcite{peffers2007design}:

\begin{enumerate}
	\item \textbf{Problem identification and motivation} --- late or missing answers to customer messages in small
	businesses, and the untrustworthiness of general-purpose chatbots (Chapter~\ref{Chapter1}).
	\item \textbf{Objectives of a solution} --- answers in Kyrgyz and Russian across messengers, no invented facts,
	bookings and orders only when valid, setup without integration work.
	\item \textbf{Design and development} --- the Joop AI platform, built iteratively.
	\item \textbf{Demonstration} --- deployment in pilot businesses in Bishkek.
	\item \textbf{Evaluation} --- scripted conversations against the real model, native-speaker review and pilot metrics.
	\item \textbf{Communication} --- this thesis, the open-source repository and its evaluation set.
\end{enumerate}

\section{System Design}

The artefact is a single Python web service (FastAPI) with a relational database (SQLite) whose schema is managed by
versioned migrations. Its main components are:

\begin{itemize}
	\item \textbf{Channel connectors} for the Telegram Bot API, the WhatsApp Cloud API, the Instagram API and a web chat.
	Incoming webhooks are verified (secret token or signature), acknowledged immediately and processed in the background;
	every message keeps its channel identifier, so a delivery repeated by the channel is processed only once.
	\item \textbf{Prompt builder}, which gives the model only the data of one business: price list, staff, working hours,
	frequent questions, free time slots and the last messages of the conversation.
	\item \textbf{Reply engine}, which asks the model for a structured JSON answer containing the reply text and,
	optionally, a proposed booking, order or hand-off to a person.
	\item \textbf{Validation layer}, which executes a proposed action only if it is valid: the slot exists, belongs to
	this business, is in the future and is still free; the ordered items exist and are in stock. A slot is claimed with a
	single conditional database update, so two customers can never receive the same time. If validation fails, the
	customer receives real alternatives instead of the model's confirmation.
	\item \textbf{Web panel} for the business: price list, staff and schedules, bookings, customers, sales funnel, cash
	and stock, broadcasts, staff accounts with roles, and a test chat in which the owner talks to the agent before
	connecting a channel.
	\item \textbf{Scheduler} for reminders, confirmations, follow-ups and the owner's morning summary.
\end{itemize}

This separation --- \enquote{the model proposes, the code decides} --- is the answer proposed for RQ1. Sector-specific
behaviour (labels, allowed outcomes such as bookings, orders or room reservations, rules and demo data) is described in
\emph{sector packs}, which is the answer proposed for RQ4; twenty packs are planned and implemented, from dental clinics to
hotels and wholesale.

\section{Data}

The study uses four kinds of data:

\begin{enumerate}
	\item \textbf{Business profiles.} For each sector pack, a realistic demo business (price list, staff, working hours,
	frequent questions) written for Bishkek. Pilot businesses provide their own real profiles.
	\item \textbf{Scripted evaluation conversations.} A set of test conversations, each with one or more customer turns
	and an expected outcome: a fact that must appear in the reply, a fact that must not be invented, a lead or order that
	must be saved, a hand-off that must happen, and the language of the reply. At the time of writing the set contains 310
	cases across a general set and the twenty sector packs, in Russian and Kyrgyz.
	\item \textbf{Pilot conversation logs.} Messages, bookings, orders and hand-offs recorded by the platform in pilot
	businesses, together with the time of each message relative to the business's working hours.
	\item \textbf{Native-speaker judgments.} Ratings of the Kyrgyz replies for correctness and naturalness by a native
	speaker of Kyrgyz.
\end{enumerate}

Personal data of customers (names and phone numbers) are used only inside the platform for the business's service; for
the thesis, results are reported in aggregate and any quoted conversation is anonymised.

\section{Steps}

\begin{enumerate}
	\item \textbf{Requirements analysis.} Study how target businesses receive messages, list the frequent questions and
	actions per sector, and review competing products.
	\item \textbf{Iterative development.} Build the platform in small increments; every change goes through a pull
	request with automated tests in continuous integration. At the time of writing the repository has more than 130
	merged pull requests and more than 500 automated tests \parencite{asilov2026joop}.
	\item \textbf{Bilingual support.} Detect the language of each message, keep Kyrgyz names of services and questions,
	and check in evaluation that Kyrgyz messages receive Kyrgyz replies (RQ3).
	\item \textbf{Offline evaluation.} Run the scripted conversations against the real language model and compute, per
	case and per sector: pass rate; rate of invented facts (a forbidden fact appears); action accuracy (a lead, booking or
	order is saved exactly when expected); hand-off accuracy; language match for Kyrgyz cases; and cost per conversation
	from the model's token usage (RQ2, RQ3).
	\item \textbf{Native-speaker review.} A native Kyrgyz speaker rates a sample of Kyrgyz replies; disagreements with the
	automatic language check are analysed.
	\item \textbf{Pilot deployment.} Connect the platform to the messengers of pilot businesses for several weeks and
	measure: share of conversations answered outside working hours; number of bookings and orders made by the agent;
	share handed to staff; time to first reply; and cost per conversation (RQ5).
	\item \textbf{Analysis and writing.} Compare results with the objectives, describe failures and their causes, and
	write the thesis.
\end{enumerate}

Prompt changes are accepted only after the offline evaluation shows that they do not reduce the pass rate, so the agent's
quality is measured before and after every change rather than judged by impression.

\section{Tools and Technologies}

Python, FastAPI, SQLModel and Alembic for the service and the database schema; Jinja2 templates for the web panel;
pytest and GitHub Actions for automated testing; the Anthropic API for the language model; the Telegram, WhatsApp and
Instagram business APIs for the channels; and a speech-to-text service for voice messages.

\section{Ethical Considerations}

The agent tells customers at the first contact that they are talking to an AI assistant, and any customer can reach a
person. Channel tokens and other secrets are stored encrypted and are never written to logs. Each business sees only its
own data. Pilot businesses take part voluntarily and may stop at any time.

\section{Time Budget}

\begin{table}[H]
\centering
\begin{tabular}{|p{6.2cm}|p{4.2cm}|p{3cm}|}
\hline
\textbf{Activity} & \textbf{Period} & \textbf{Status} \\
\hline
Topic selection and approval & October 2026 & Done \\
\hline
Requirements analysis and competitor review & October 2026 & Done \\
\hline
Architecture and core platform (channels, engine, validation, panel) & October -- November 2026 & Done \\
\hline
Sector packs and bilingual support & October -- November 2026 & Done \\
\hline
Thesis proposal & November 2026 & In progress \\
\hline
Offline evaluation on scripted conversations & November 2026 & Planned \\
\hline
Native-speaker review of Kyrgyz replies & November 2026 & Planned \\
\hline
Pilot deployment in businesses & November -- December 2026 & Planned \\
\hline
Analysis and writing of the thesis & November -- December 2026 & Planned \\
\hline
Submission of the thesis & by 1 December 2026 & Planned \\
\hline
\end{tabular}
\caption{Time budget of the study}
\label{tab:timebudget}
\end{table}

\section{Summary}

The study will design, build and evaluate Joop AI following design science research. The platform separates
conversation (by the language model) from business actions (validated in code), supports Kyrgyz and Russian across four
channels, and is generalised through sector packs. It will be evaluated with explicit, automatically checked scripted
conversations, a native-speaker review of Kyrgyz replies and pilot deployments that measure customers served outside
working hours, bookings made by the agent and the cost of each conversation.
